\section{Описание}

Требуется реализовать поиск поступающих на вход образцов (паттернов) в заранее известном тексте. Для эффективного решения задачи применяется предварительная индексация текста — построение \textbf{суффиксного массива}.

Для построения суффиксного массива напрямую используется алгоритм \textbf{Манбера-Майерса}. Основная идея алгоритма заключается в итеративной сортировке префиксов суффиксов, длины которых являются степенями двойки ($1, 2, 4, 8 \dots$). На каждой итерации длины сравниваемых префиксов удваиваются, что позволяет достичь временной сложности $O(N \log N)$. Для достижения такой оценки на каждом шаге применяется сортировка подсчётом (стабильная сортировка) по классам эквивалентности.

Поиск паттернов в построенном массиве осуществляется с помощью \textbf{бинарного поиска}. Поскольку суффиксный массив содержит лексикографически отсортированные суффиксы текста, все вхождения искомого паттерна будут располагаться в массиве непрерывным блоком. С помощью модифицированных функций верхней и нижней границы (\texttt{lower\_bound} и \texttt{upper\_bound}) алгоритм за логарифмическое время находит границы этого блока.

Для корректной лексикографической сортировки и обработки краевых случаев к исходному тексту добавляется специальный терминальный символ \texttt{\$}, который лексикографически меньше любого другого символа алфавита.

\pagebreak

\section{Исходный код}

Реализация состоит из функции построения суффиксного массива и основного цикла программы, отвечающего за потоковое чтение паттернов и их поиск.

\subsection{Алгоритм Манбера-Майерса}

Функция \texttt{build\_sa} отвечает за построение суффиксного массива.
\begin{itemize}
    \item \texttt{p} — массив перестановок (сам суффиксный массив), хранящий индексы начала суффиксов.
    \item \texttt{c} — массив классов эквивалентности. Если два суффикса совпадают на текущей длине префикса, они имеют одинаковый класс.
    \item \texttt{cnt} — массив счётчиков для линейной сортировки подсчётом.
\end{itemize}
Алгоритм начинает с сортировки циклических сдвигов длины 1, после чего в цикле удваивает длину префикса ($1 \ll h$), пересчитывая классы эквивалентности до тех пор, пока количество уникальных классов не станет равным длине текста.

\subsection{Бинарный поиск паттернов}

Для поиска паттернов в функции \texttt{main} используются лямбда-выражения \texttt{comp\_low} и \texttt{comp\_up}. Они передаются в стандартные алгоритмы \texttt{std::lower\_bound} и \texttt{std::upper\_bound} для бинарного поиска начала и конца отрезка в суффиксном массиве, в котором суффиксы начинаются с искомого паттерна. Метод \texttt{text.compare} позволяет сравнивать только нужный префикс суффикса длиной \texttt{m}.

\begin{table}[!ht]
\centering
\begin{tabular}{|p{7.5cm}|p{7.5cm}|}
\hline
\rowcolor{lightgray}
\multicolumn{2}{|c|}{main.cpp} \\
\hline
\texttt{void build\_sa(const std::string\& s, std::vector<int>\& p)} & Функция построения суффиксного массива алгоритмом Манбера-Майерса. \\
\hline
\rowcolor{lightgray}
\multicolumn{2}{|c|}{main.cpp} \\
\hline
\texttt{comp\_low} / \texttt{comp\_up} & Лямбда-компараторы для нахождения границ диапазона совпадений через бинарный поиск. \\
\hline
\rowcolor{lightgray}
\multicolumn{2}{|c|}{main.cpp} \\
\hline
\texttt{std::vector<int> occurrences} & Вектор для хранения позиций найденного паттерна с последующей сортировкой для форматированного вывода. \\
\hline
\end{tabular}
\end{table}

\pagebreak

\begin{lstlisting}[language=C++]
#include <iostream>
#include <vector>
#include <string>
#include <algorithm>

void build_sa(const std::string& s, std::vector<int>& p) {
    int n = s.length();
    const int alphabet = 256;
    p.assign(n, 0);
    std::vector<int> c(n, 0), cnt(std::max(alphabet, n), 0);

    for (int i = 0; i < n; i++) cnt[(unsigned char)s[i]]++;
    for (int i = 1; i < alphabet; i++) cnt[i] += cnt[i - 1];
    for (int i = 0; i < n; i++) p[--cnt[(unsigned char)s[i]]] = i;

    c[p[0]] = 0;
    int classes = 1;
    for (int i = 1; i < n; i++) {
        if (s[p[i]] != s[p[i - 1]]) classes++;
        c[p[i]] = classes - 1;
    }

    std::vector<int> pn(n), cn(n);
    for (int h = 0; (1 << h) < n; ++h) {
        for (int i = 0; i < n; i++) {
            pn[i] = p[i] - (1 << h);
            if (pn[i] < 0) pn[i] += n;
        }
        
        fill(cnt.begin(), cnt.begin() + classes, 0);
        for (int i = 0; i < n; i++) cnt[c[pn[i]]]++;
        for (int i = 1; i < classes; i++) cnt[i] += cnt[i - 1];
        for (int i = n - 1; i >= 0; i--) p[--cnt[c[pn[i]]]] = pn[i];

        cn[p[0]] = 0;
        classes = 1;

        for (int i = 1; i < n; i++) {
            int mid1 = p[i] + (1 << h);
            if (mid1 >= n) mid1 -= n;
            
            int mid2 = p[i - 1] + (1 << h);
            if (mid2 >= n) mid2 -= n;

            if (c[p[i]] != c[p[i - 1]] || c[mid1] != c[mid2]) {
                ++classes;
            }
            cn[p[i]] = classes - 1;
        }
        c = cn;

        if (classes == n) break;
    }
}

int main() {
    std::ios_base::sync_with_stdio(false);
    std::cin.tie(NULL);

    std::string text;
    if (!(std::cin >> text)) return 0;
    
    text += '$';

    std::vector<int> p;
    build_sa(text, p);

    std::string pattern;
    int pattern_idx = 1;

    std::vector<int> occurrences;
    
    while (std::cin >> pattern) {
        int m = pattern.length();

        auto comp_low = [&](int suffix_idx, const std::string& pat) {
            return text.compare(suffix_idx, m, pat) < 0;
        };

        auto comp_up = [&](const std::string& pat, int suffix_idx) {
            return text.compare(suffix_idx, m, pat) > 0;
        };

        auto it_start = lower_bound(p.begin(), p.end(), pattern, comp_low);
        auto it_end = upper_bound(it_start, p.end(), pattern, comp_up);

        if (it_start != it_end) {
            occurrences.clear();
            
            for (auto it = it_start; it != it_end; ++it) {
                occurrences.push_back(*it + 1);
            }
            
            std::sort(occurrences.begin(), occurrences.end());

            std::cout << pattern_idx << ": ";
            for (size_t i = 0; i < occurrences.size(); ++i) {
                std::cout << occurrences[i] << (i + 1 == occurrences.size() ? "" : ", ");
            }
            std::cout << "\n";
        }
        pattern_idx++;
    }

    return 0;
}
\end{lstlisting}

\subsection{Анализ сложности}

Пусть $N$ — длина текста, $M$ — длина текущего искомого паттерна.
\begin{itemize}
    \item \textbf{Временная сложность предварительной обработки:} $O(N \log N)$. Алгоритм Манбера-Майерса требует $\log N$ итераций, на каждой из которых выполняется сортировка подсчётом за $O(N)$.
    \item \textbf{Временная сложность поиска одного паттерна:} $O(M \log N)$. Бинарный поиск требует $O(\log N)$ шагов, на каждом из которых выполняется лексикографическое сравнение строк длины до $M$.
    \item \textbf{Пространственная сложность:} $O(N)$. Суффиксный массив и вспомогательные структуры для построения (массивы $p, c, cnt, pn, cn$) требуют линейного количества памяти относительно длины текста.
\end{itemize}

Данный алгоритм является крайне эффективным для задач с однократным чтением объемного текста и последующим множественным поиском образцов, так как основное время тратится на препроцессинг, а после для каждого запроса требуется в разы меньше времени.

\pagebreak

\section{Консоль}

\begin{alltt}
root@workspaces:/workspaces/Discrete-Analysis-MAI# cmake -B build
root@workspaces:/workspaces/Discrete-Analysis-MAI# cmake --build build
root@workspaces:/workspaces/Discrete-Analysis-MAI# ./build/lab5/lab5

abcdabc
abcd
bcd
bc
1: 1
2: 2
3: 2, 6

\end{alltt}

\pagebreak